% GENERATED FILE. Edit canonical lessons, then run npm run generate:books.
% canonical-source-hash: fnv1a64:2c750469f14fbe0b
% canonical-lessons: ML-C395-purvikan, ML-C395-pingami, ML-C395-acaram, ML-C395-utbhavam, ML-C395-muttassikkatha

\chapter{Ancestors and Customs}
\label{ch:ml-ancestors-and-customs}
\hlchaptermodality{395}

\begin{chapteropening}
\textbf{By the end of this chapter:} \emph{I can say an ancestor, a successor, a custom, an origin and a grandmother's tale.}

Complete the last lesson of chapter 395: I can say an ancestor, a successor, a custom, an origin and a grandmother's tale.
\end{chapteropening}

\section[\texorpdfstring{\ml{പൂർവികൻ} (pūrvikan)}{പൂർവികൻ (pūrvikan)}]{\ml{പൂർവികൻ} (pūrvikan) — an ancestor}

\label{lesson:ML-C395-purvikan}



\begin{quote}
Before the new one: say the Malayalam for a widow, then the Malayalam for heritage.
\end{quote}

\subsection*{You'll want to know: \ml{പൂർവികൻ}}
\textbf{\ml{പൂർവികൻ}} — \emph{pūrvikan} — \textquotedblleft{}an ancestor\textquotedblright{}.

\begin{quote}
\ml{ഞങ്ങളുടെ} \ml{പൂർവികർ} \ml{കർഷകർ} \ml{ആയിരുന്നു}.
\end{quote}

\begin{quote}
\emph{ñaṅṅaḷuṭe pūrvikar karṣakar āyirunnu.}
\end{quote}

\begin{quote}
Our ancestors were farmers.
\end{quote}

Plural: \ml{പൂർവികർ}, \emph{pūrvikar}.

Usually in the plural, \ml{പൂർവികർ}, our forefathers.

\begin{cousinweb}
\textbf{\ml{പൂർവികൻ}} (\emph{pūrvikan}) is borrowed from Sanskrit \emph{pūrvika}, \textquotedblleft{}former, ancient\textquotedblright{} (Monier-Williams), from \emph{pūrva}, \textquotedblleft{}before, earlier\textquotedblright{}, with the Malayalam masculine ending -\ml{ൻ}.
\end{cousinweb}

\begin{grammarlens}[title={What you've built}]
One more word for telling what happened.
\end{grammarlens}

\subsection*{Your turn}
\begin{itemize}
\raggedright
  \item \emph{Say it:} \emph{pūrvikan}
  \item \emph{Say it:} \emph{pūrvikan}, once more
  \item \emph{Say it:} \emph{ñaṅṅaḷuṭe pūrvikar karṣakar āyirunnu.} — Our ancestors were farmers.
  \item \emph{Recall:} say the Malayalam for a widow, then the Malayalam for heritage, then say \emph{pūrvikan} again
\end{itemize}

\begin{tcolorbox}[breakable,colback=teal!4,colframe=teal!35!black,title={Before you move on}]
What does \ml{പൂർവികൻ} mean? (\textquotedblleft{}An ancestor\textquotedblright{}.) Say it once more.
\end{tcolorbox}
\section[\texorpdfstring{\ml{പിൻഗാമി} (pingāmi)}{പിൻഗാമി (pingāmi)}]{\ml{പിൻഗാമി} (pingāmi) — a successor}

\label{lesson:ML-C395-pingami}



\begin{quote}
Before the new one: say the Malayalam for heritage, then the Malayalam for an ancestor.
\end{quote}

\subsection*{You'll want to know: \ml{പിൻഗാമി}}
\textbf{\ml{പിൻഗാമി}} — \emph{pingāmi} — \textquotedblleft{}a successor\textquotedblright{}.

\begin{quote}
\ml{മാനേജർ} \ml{തന്റെ} \ml{പിൻഗാമിയെ} \ml{പരിചയപ്പെടുത്തി}.
\end{quote}

\begin{quote}
\emph{mānējar tanṟe pingāmiye paricayappeṭutti.}
\end{quote}

\begin{quote}
The manager introduced his successor.
\end{quote}

\ml{പിൻ}, after, + -\ml{ഗാമി}, one who goes: the one who comes after.

\begin{cousinweb}
\textbf{\ml{പിൻഗാമി}} (\emph{pingāmi}) joins the native \ml{പിൻ} (\emph{pin}), \textquotedblleft{}behind, after\textquotedblright{} (DEDR 4205), and Sanskrit \emph{gāmin}, \textquotedblleft{}going\textquotedblright{} (CDIAL 4134): the one who goes after.
\end{cousinweb}

\begin{grammarlens}[title={What you've built}]
One more word for telling what happened.
\end{grammarlens}

\subsection*{Your turn}
\begin{itemize}
\raggedright
  \item \emph{Say it:} \emph{pingāmi}
  \item \emph{Say it:} \emph{pingāmi}, once more
  \item \emph{Say it:} \emph{mānējar tanṟe pingāmiye paricayappeṭutti.} — The manager introduced his successor.
  \item \emph{Recall:} say the Malayalam for heritage, then the Malayalam for an ancestor, then say \emph{pingāmi} again
\end{itemize}

\begin{tcolorbox}[breakable,colback=teal!4,colframe=teal!35!black,title={Before you move on}]
What does \ml{പിൻഗാമി} mean? (\textquotedblleft{}A successor\textquotedblright{}.) Say it once more.
\end{tcolorbox}
\section[\texorpdfstring{\ml{ആചാരം} (ācāraṁ)}{ആചാരം (ācāraṁ)}]{\ml{ആചാരം} (ācāraṁ) — a custom; a rite}

\label{lesson:ML-C395-acaram}



\begin{quote}
Before the new one: say the Malayalam for an ancestor, then the Malayalam for a successor.
\end{quote}

\subsection*{You'll want to know: \ml{ആചാരം}}
\textbf{\ml{ആചാരം}} — \emph{ācāraṁ} — \textquotedblleft{}a custom; a rite\textquotedblright{}.

\begin{quote}
\ml{ഓരോ} \ml{നാടിനും} \ml{അതിന്റേതായ} \ml{ആചാരങ്ങൾ} \ml{ഉണ്ട്}.
\end{quote}

\begin{quote}
\emph{ōrō nāṭinuṁ atinṟētāya ācāraṅṅaḷ uṇṭŭ.}
\end{quote}

\begin{quote}
Every land has its own customs.
\end{quote}

Plural: \ml{ആചാരങ്ങൾ}, \emph{ācāraṅṅaḷ}.

Plural \ml{ആചാരങ്ങൾ}, the customs and rites of a place.

\begin{cousinweb}
\textbf{\ml{ആചാരം}} (\emph{ācāraṁ}) is borrowed from Sanskrit \emph{ācāra}, \textquotedblleft{}mode of action, conduct; custom, usage\textquotedblright{} (CDIAL 1071).
\end{cousinweb}

\begin{grammarlens}[title={What you've built}]
One more word for telling what happened.
\end{grammarlens}

\subsection*{Your turn}
\begin{itemize}
\raggedright
  \item \emph{Say it:} \emph{ācāraṁ}
  \item \emph{Say it:} \emph{ācāraṁ}, once more
  \item \emph{Say it:} \emph{ōrō nāṭinuṁ atinṟētāya ācāraṅṅaḷ uṇṭŭ.} — Every land has its own customs.
  \item \emph{Recall:} say the Malayalam for an ancestor, then the Malayalam for a successor, then say \emph{ācāraṁ} again
\end{itemize}

\begin{tcolorbox}[breakable,colback=teal!4,colframe=teal!35!black,title={Before you move on}]
What does \ml{ആചാരം} mean? (\textquotedblleft{}A custom; a rite\textquotedblright{}.) Say it once more.
\end{tcolorbox}
\section[\texorpdfstring{\ml{ഉത്ഭവം} (utbhavaṁ)}{ഉത്ഭവം (utbhavaṁ)}]{\ml{ഉത്ഭവം} (utbhavaṁ) — an origin; a source}

\label{lesson:ML-C395-utbhavam}



\begin{quote}
Before the new one: say the Malayalam for a successor, then the Malayalam for a custom.
\end{quote}

\subsection*{You'll want to know: \ml{ഉത്ഭവം}}
\textbf{\ml{ഉത്ഭവം}} — \emph{utbhavaṁ} — \textquotedblleft{}an origin; a source\textquotedblright{}.

\begin{quote}
\ml{ആ} \ml{പുഴയുടെ} \ml{ഉത്ഭവം} \ml{ഒരു} \ml{മലയിൽ} \ml{ആണ്}.
\end{quote}

\begin{quote}
\emph{ā puḻayuṭe utbhavaṁ oru malayil āṇŭ.}
\end{quote}

\begin{quote}
The source of that river is on a hill.
\end{quote}

Also the origin of a word or a custom.

\begin{cousinweb}
\textbf{\ml{ഉത്ഭവം}} (\emph{utbhavaṁ}) is borrowed from Sanskrit \emph{udbhava}, \textquotedblleft{}origin, birth\textquotedblright{} (CDIAL 2039), from \emph{ud-bhū}, \textquotedblleft{}to come up, to arise\textquotedblright{}.
\end{cousinweb}

\begin{grammarlens}[title={What you've built}]
One more word for telling what happened.
\end{grammarlens}

\subsection*{Your turn}
\begin{itemize}
\raggedright
  \item \emph{Say it:} \emph{utbhavaṁ}
  \item \emph{Say it:} \emph{utbhavaṁ}, once more
  \item \emph{Say it:} \emph{ā puḻayuṭe utbhavaṁ oru malayil āṇŭ.} — The source of that river is on a hill.
  \item \emph{Recall:} say the Malayalam for a successor, then the Malayalam for a custom, then say \emph{utbhavaṁ} again
\end{itemize}

\begin{tcolorbox}[breakable,colback=teal!4,colframe=teal!35!black,title={Before you move on}]
What does \ml{ഉത്ഭവം} mean? (\textquotedblleft{}An origin; a source\textquotedblright{}.) Say it once more.
\end{tcolorbox}
\section[\texorpdfstring{\ml{മുത്തശ്ശിക്കഥ} (muttaśśikkatha)}{മുത്തശ്ശിക്കഥ (muttaśśikkatha)}]{\ml{മുത്തശ്ശിക്കഥ} (muttaśśikkatha) — a grandmother's tale}

\label{lesson:ML-C395-muttassikkatha}



\begin{quote}
Before the new one: say the Malayalam for a custom, then the Malayalam for an origin.
\end{quote}

\subsection*{You'll want to know: \ml{മുത്തശ്ശിക്കഥ}}
\textbf{\ml{മുത്തശ്ശിക്കഥ}} — \emph{muttaśśikkatha} — \textquotedblleft{}a grandmother's tale\textquotedblright{}.

\begin{quote}
\ml{രാത്രി} \ml{മുത്തശ്ശി} \ml{ഞങ്ങൾക്ക്} \ml{ഒരു} \ml{മുത്തശ്ശിക്കഥ} \ml{പറഞ്ഞുതന്നു}.
\end{quote}

\begin{quote}
\emph{rātri muttaśśi ñaṅṅaḷkkŭ oru muttaśśikkatha paṟaññutannu.}
\end{quote}

\begin{quote}
At night grandmother told us an old tale.
\end{quote}

A fixed word, not two: also said of a story nobody believes.

\begin{cousinweb}
\textbf{\ml{മുത്തശ്ശിക്കഥ}} (\emph{muttaśśikkatha}) joins \ml{മുത്തശ്ശി} (\emph{muttaśśi}), \textquotedblleft{}grandmother\textquotedblright{}, and \ml{കഥ} (\emph{katha}), \textquotedblleft{}a story\textquotedblright{}, borrowed from Sanskrit \emph{kathā} (CDIAL 2704).
\end{cousinweb}

\begin{grammarlens}[title={What you've built}]
One more word for telling what happened.
\end{grammarlens}

\subsection*{Your turn}
\begin{itemize}
\raggedright
  \item \emph{Say it:} \emph{muttaśśikkatha}
  \item \emph{Say it:} \emph{muttaśśikkatha}, once more
  \item \emph{Say it:} \emph{rātri muttaśśi ñaṅṅaḷkkŭ oru muttaśśikkatha paṟaññutannu.} — At night grandmother told us an old tale.
  \item \emph{Recall:} say the Malayalam for a custom, then the Malayalam for an origin, then say \emph{muttaśśikkatha} again
\end{itemize}

\begin{tcolorbox}[breakable,colback=teal!4,colframe=teal!35!black,title={Before you move on}]
What does \ml{മുത്തശ്ശിക്കഥ} mean? (\textquotedblleft{}A grandmother's tale\textquotedblright{}.) Say it once more.
\end{tcolorbox}
